\chapter{Warts}

\section*{Definition}
Warts (verrucae) are benign epidermal proliferations caused by human papillomavirus (HPV) infection of keratinocytes. They are classified by morphology and location: common (verruca vulgaris), plantar (verruca plantaris), flat (verruca plana), filiform, and genital (condyloma acuminatum).

\section*{Pathophysiology}
HPV infects basal keratinocytes through microabrasions in the stratum corneum, where viral replication induces hyperkeratosis, acanthosis, and papillomatosis. The virus remains intraepithelial, evading immune detection through downregulation of interferon signaling and MHC class I expression. Cell-mediated immunity ultimately clears infection, explaining the higher prevalence of warts in immunosuppressed individuals and their tendency to regress spontaneously in immunocompetent hosts.

\section*{Etiology}
\begin{itemize}
  \item \textbf{Viral:} HPV types 1, 2, 4, 27, 57 (common/plantar), types 3, 10 (flat), types 6, 11, 16, 18 (genital)
  \item \textbf{Transmission:} Direct skin-to-skin contact, autoinoculation, fomites (shared towels, gym floors, swimming pools)
  \item \textbf{Risk factors:} Children and young adults, immunosuppression (HIV, transplant, chemotherapy), chronic skin conditions causing barrier disruption, occupations with frequent hand wetting (butchers, fish handlers)
\end{itemize}

\section*{Indications for Evaluation}
Seek care if:
\begin{itemize}
  \item Painful warts (especially plantar) interfering with walking or daily activities
  \item Rapidly spreading, numerous, or enlarging warts, particularly in immunocompromised patients
  \item Atypical appearance: pigmented, ulcerated, or bleeding lesion (rule out squamous cell carcinoma)
  \item Genital or perianal warts (sexual abuse concern in children; STI evaluation in adults)
\end{itemize}

\section*{Related Symptoms}
\begin{itemize}
  \item Molluscum contagiosum, seborrheic keratosis, corns and calluses, acrochordon, squamous cell carcinoma
\end{itemize}
\textbf{Note:} Up to 65\% of common warts in children resolve spontaneously within 2 years, but treatment is often pursued due to pain, cosmetic concerns, or to reduce transmission.

\section*{Self-Care / Home Management}
\begin{itemize}
  \item Over-the-counter salicylic acid 17--40\% topical solution or patch applied daily after soaking and gentle debridement
  \item Occlusion with duct tape (cut to size, left on 6 days, removed, soaked, debrided, reapplied for 8 weeks)
  \item File down plantar warts with a pumice stone between treatments; avoid sharing nail clippers or files
\end{itemize}

\section*{Diagnostic Approach}
\begin{itemize}
  \item Clinical diagnosis: Verrucous surface, loss of skin lines, pinpoint black dots (thrombosed capillaries), pain with lateral compression (plantar warts)
  \item Dermatoscopy: Thrombosed capillaries, papillomatosis, and interrupted dermatoglyphics
  \item Biopsy if atypical features (rapid growth, pigmentation, ulceration, immunosuppression with suspicion of malignancy)
\end{itemize}

\section*{Treatment / Management Overview}
\begin{itemize}
  \item First-line: Cryotherapy with liquid nitrogen (10--30 sec freeze, 1--3 sessions at 2--4 week intervals); salicylic acid as self-administered alternative
  \item Second-line: Cantharidin (blistering agent), intralesional Candida/MMP antigen (immunotherapy), pulsed-dye laser, photodynamic therapy
  \item Refractory: Bleomycin intralesional injection, topical imiquimod 5\% cream, topical 5-fluorouracil, or surgical excision (last resort due to scarring)
  \item Genital warts: Patient-applied podofilox or imiquimod; provider-administered cryotherapy, TCA, or electrosurgery
\end{itemize}

\clearpage
